\chapter{Geometric and Equivariant Deep Learning}
\label{ch:equivariant}

\epigraph{\itshape Symmetry, as wide or as narrow as you may define its
meaning, is one idea by which man through the ages has tried to
comprehend and create order, beauty, and perfection.}{Hermann Weyl}

\section{Why symmetry matters}
Physical systems exhibit symmetries: translations, rotations,
reflections, Galilean boosts, gauge transformations. A model that
respects these symmetries by construction (\emph{equivariant}) requires
less data, generalises better, and respects conservation laws --- a
direct consequence of Noether's theorem.

\begin{defn}{Equivariance}{equiv}
A map $\Phi:\X\to\Y$ is \emph{equivariant} with respect to a group $G$
acting on $\X$ via $\rho_X$ and on $\Y$ via $\rho_Y$ if
$\Phi(\rho_X(g)\bm{x}) = \rho_Y(g)\Phi(\bm{x})$ for all $g\in G$.
Invariance is the special case $\rho_Y\equiv \mathrm{id}$.
\end{defn}

\section{Group-equivariant CNNs}
G-CNNs generalise convolution: weights are functions on $G$ rather than
$\R^d$, and convolution becomes integration over the group. Cohen \&
Welling (2016) introduced the framework for finite groups
($\mathbb{Z}_4$, dihedral). Steerable CNNs use representation theory to
construct continuous-group equivariant filters.

\section{$E(3)$ and $SE(3)$ equivariance}
For 3D physics, $E(3)$ (translations + rotations + reflections) is the
natural symmetry group. Two leading paradigms:

\paragraph{Tensor field networks / SE(3)-Transformer / NequIP.}
Activations are sums of irreducible representations
$\mathbf{V}=\bigoplus_\ell \mathbf{V}^{(\ell)}$ of $SO(3)$. Layers
combine these via tensor products and Clebsch--Gordan
coefficients~\cite{thomas2018tfn,fuchs2020se3,batzner2022nequip}.

\paragraph{E(n)-equivariant graph neural networks (EGNN).}
Maintain scalar features $\bm{h}_i$ and equivariant coordinates
$\bm{x}_i$; updates use only invariants like $\|\bm{x}_i-\bm{x}_j\|$.
Lightweight and effective~\cite{satorras2021egnn}.

\paragraph{Geometric algebra networks.}
GATr (Brehmer et al.) operate in the projective geometric algebra of
3D Euclidean space, encoding rotations, translations, and reflections
naturally.

\section{N-body dynamics}
Equivariant networks dominate learned $n$-body simulators:
\begin{itemize}
\item \textbf{Molecular dynamics}: NequIP, Allegro, MACE, SO3krates,
GemNet --- machine-learning force fields with $SO(3)$ equivariance.
\item \textbf{Astrophysics}: equivariant graph networks for galaxy
mergers, halo finders, cosmological parameter inference. Particularly
relevant when conserving angular momentum and energy is critical for
long-horizon stability.
\item \textbf{Particle physics}: ParticleNet, LorentzNet (Lorentz
equivariance) for jet tagging and event classification.
\end{itemize}

\begin{insight}{Equivariance as inductive bias}
The single biggest lesson of geometric deep learning: \emph{the right
inductive bias outweighs more data}. NequIP/Allegro reach DFT accuracy
on chemistry tasks with $\sim$10$\times$ less data than non-equivariant
counterparts.
\end{insight}

\section{Hamiltonian / Lagrangian + equivariance}
Combining mechanical priors with symmetries gives the strongest
extrapolation. Examples:
\begin{itemize}
\item Constrained Hamiltonian neural networks for spherical pendulums.
\item Lagrangian neural networks with $SE(3)$ equivariance for
articulated robotics.
\item Symplectic + equivariant ODE-Nets for n-body simulations.
\end{itemize}

\section{Mesh-based and graph-based PDE surrogates}
\paragraph{MeshGraphNets~\cite{pfaff2021mgn}}
Learn dynamics on mesh graphs with edge-conditioned message passing;
flexible across geometries; foundational for many CFD, soft-body, and
cloth simulators.
\paragraph{Lagrangian particle GNNs~\cite{sanchez2020graphnet}.}
Treat particles as nodes; predict accelerations; integrate explicitly.
Used for fluids, granular media.

\section{Open challenges}
\begin{itemize}
\item Computational cost of high-$\ell$ tensor products (mitigated by
MACE / SO3krates / Cartesian-tensor reformulations).
\item Equivariance under \emph{approximate} symmetries (e.g.\ slight
gauge breaking).
\item Combining symmetry with geometric heterogeneity (boundaries,
defects).
\item Foundation models that are simultaneously equivariant and large.
\end{itemize}

\section{Recipe: choosing an equivariant model}
\begin{recipe}{Equivariant architecture selection}
\begin{enumerate}
\item Symmetry group of your problem (E(3)? SE(3)? Lorentz? gauge?).
\item Required output type (scalar invariant? vector? tensor?).
\item Cost budget: EGNN $<$ MACE $<$ TFN/SE(3)-T.
\item Conservation requirements: pair with HNN/LNN if energy/momentum
conservation is crucial.
\end{enumerate}
\end{recipe}
